\hypertarget{incremental}{}
\hypertarget{the-incremental-test-dynamic-forecast-or-stable-ordering}{%
\chapter{The incremental test: dynamic forecast or stable
ordering?}\label{the-incremental-test-dynamic-forecast-or-stable-ordering}}

\hypertarget{the-identification-problem}{%
\section{The identification problem}\label{the-identification-problem}}

Suppose FF3 persistently underprices a \emph{fixed} set of characteristic
cells --- small-growth structurally poor, small-value strong. Formally, if
both the training score and the future alpha are noisy readings of the same
fixed cell effect $c_i$ (\S\ref{rq-difficulties}),
\begin{equation}\label{eq:ch12-fixedeffect}
\hat s_{it} = c_i + u_{it},\qquad A^{+}_{i,t} = c_i + w_{i,t},
\end{equation}
then their cross-sectional correlation is positive through the shared $c_i$
alone, with \emph{no} time-varying information in the ranking. Past alpha
therefore predicts future alpha even when annual re-estimation adds nothing.
The forward \emph{level} test of Chapter~11 cannot distinguish this stable
ordering from genuine dynamic ranking skill; separating the two requires a
test of \emph{equal predictive accuracy} against a fixed comparator
(\S\ref{prelim-forecast}).

\hypertarget{comparator-rankings}{%
\section{Comparator rankings}\label{comparator-rankings}}

\begin{longtable}[]{@{}p{0.20\textwidth}p{0.34\textwidth}p{0.36\textwidth}@{}}
\toprule\noalign{}
Comparator & Uses future data? & Purpose \\
\midrule\noalign{}
\endhead
\bottomrule\noalign{}
\endlastfoot
Size-ascending / B/M-descending gradient & No & Naive characteristic
ordering \\
Diagonal value-minus-size gradient & No & Alternative simple
characteristic ordering \\
Expanding-window ranking & No & Fair strong comparator using all data
available at each cutoff \\
Full-sample alpha oracle & \textbf{Yes} & Look-ahead upper bound, never
a tradable strategy \\
\end{longtable}

\hypertarget{increment-definition}{%
\section{Increment definition}\label{increment-definition}}

For each evaluation window $w$, let $\rho^{\mathrm{roll}}_w$ and
$\rho^{\mathrm{bench}}_w$ be the predictive scores --- rank correlation, or
top-minus-bottom quintile spread --- of the rolling ranking and a comparator
on the \emph{same} future observations \eqref{eq:ch11-fwdalpha}, and form
the per-window skill differential
\begin{equation}\label{eq:ch12-increment}
d_w = \rho^{\mathrm{roll}}_w - \rho^{\mathrm{bench}}_w .
\end{equation}
Its mean $\bar d$ standardised by a HAC standard error is precisely the
\emph{Diebold--Mariano} statistic of equal predictive accuracy
\eqref{eq:dm}, tested one-sided against the prespecified alternative
$\mathbb{E}(d_w)>0$; because the comparators are re-estimated on
rolling/expanding schemes it carries the Giacomini--White interpretation
(\S\ref{prelim-forecast}). A block bootstrap corroborates the HAC $t$-test.
The fair verdict rests on the expanding-window comparator, because both
rankings then use only information available at the cutoff.

\hypertarget{a-numerical-example-that-separates-level-from-increment}{%
\section{A numerical example that separates level from
increment}\label{a-numerical-example-that-separates-level-from-increment}}

Suppose four evaluation years produce rolling-rank correlations (0.12,
0.08, 0.16, 0.04) and expanding-rank correlations (0.15, 0.11, 0.18,
0.07). Both methods predict positively:

\begin{equation}\label{eq:ch12-means}
\overline{\rho}_{\mathrm{rolling}} = 0.10,
\qquad
\overline{\rho}_{\mathrm{expanding}} = 0.1275 .
\end{equation}

But the paired increments are $(-0.03,-0.03,-0.02,-0.03)$, with mean
$\bar d=-0.0275$. Thus a positive forward level is perfectly compatible with
a \emph{negative} incremental contribution from rolling re-estimation: the
shared fixed cross-sectional signal $c_i$ explains the positive level, while
discarding older observations to re-estimate on a 120-month window adds
estimation noise. Only the paired differential --- not either level ---
answers the dynamic-skill question.

\hypertarget{interpretation-of-the-reported-evidence}{%
\section{Interpretation of the reported
evidence}\label{interpretation-of-the-reported-evidence}}

\begin{itemize}
\tightlist
\item
  The rolling rank beats the two naive characteristic gradients.
\item
  It does not beat the expanding-window rank: mean correlation increment
  \ensuremath{-}0.042, one-sided p = 0.97. For the prespecified alternative
  E(\ensuremath{\Delta })\textgreater0, such a large p-value means the estimate lies in the
  wrong direction; it is not evidence that the two procedures are
  equivalent.
\item
  It is dominated by the full-sample oracle, as expected for an upper
  bound with look-ahead.
\end{itemize}

\textbf{Central conclusion.} The out-of-sample relation is consistent
with a stable, irregular cross-sectional pattern of benchmark-relative
residual returns. Calling this ``mispricing'' would require additional
economic identification. Discarding older data and re-estimating from a
120-month rolling window adds noise rather than dynamic forecasting
information.

\hypertarget{minimal-implementation-pattern}{%
\section{Minimal implementation
pattern}\label{minimal-implementation-pattern}}

\begin{verbatim}
for cutoff in cutoffs:
    future = future_alpha(cutoff, horizon=12, frozen_betas=True)
    rolling = rank_from_last_120_months(cutoff)
    expanding = rank_from_all_history_to(cutoff)

    rho_roll = spearmanr(rolling, future).statistic
    rho_expand = spearmanr(expanding, future).statistic
    increments.append(rho_roll - rho_expand)

estimate, se, p_one_sided = hac_mean_test(increments, alternative="greater")
\end{verbatim}

